\documentclass[10pt,a4paper]{article}

% =========================================================
% ENCODING AND PDF TEXT EXTRACTION
% Compile using pdfLaTeX.
% =========================================================
\usepackage[T1]{fontenc}
\usepackage[utf8]{inputenc}
\usepackage{lmodern}
\input{glyphtounicode}
\pdfgentounicode=1

% =========================================================
% PAGE LAYOUT
% =========================================================
\usepackage[margin=0.65in]{geometry}
\usepackage{microtype}
\usepackage{enumitem}
\usepackage{titlesec}
\usepackage{needspace}
\usepackage[hidelinks]{hyperref}

\renewcommand{\familydefault}{\sfdefault}
\pagestyle{empty}
\raggedright
\raggedbottom

\setlength{\parindent}{0pt}
\setlength{\parskip}{2pt}
\setlength{\emergencystretch}{2em}
\urlstyle{same}

\hypersetup{
  pdftitle={Apu Das - Micron Scholarship Resume},
  pdfauthor={Apu Das},
  pdfsubject={Semiconductor Memory and Device Reliability}
}

% =========================================================
% SECTION AND LIST FORMATTING
% =========================================================
\titleformat{\section}
  {\normalsize\bfseries}
  {}{0pt}{}
  [\titlerule]

\titlespacing*{\section}{0pt}{9pt}{4pt}

\setlist[itemize]{
  leftmargin=1.25em,
  itemsep=2pt,
  topsep=2pt,
  parsep=0pt,
  partopsep=0pt
}

\setlist[enumerate]{
  leftmargin=1.5em,
  itemsep=6pt,
  topsep=3pt,
  parsep=0pt,
  partopsep=0pt
}

% Single-column reading order, without tables or text boxes.
\newcommand{\entry}[3]{%
  \par\addvspace{4pt}%
  \needspace{4\baselineskip}%
  \textbf{#1}\par
  #2\par
  \textit{#3}\par
}

\begin{document}

% =========================================================
% HEADER
% =========================================================
{\fontsize{24}{27}\selectfont\bfseries APU DAS}\par
\textbf{PhD Researcher | Semiconductor Memory and Device Reliability}\par
College of Semiconductor Research, National Tsing Hua University\par
Hsinchu, Taiwan |
\href{mailto:apu.das@gapp.nthu.edu.tw}{apu.das@gapp.nthu.edu.tw}
| +886 910 083 544\par
LinkedIn:
\href{https://linkedin.com/in/apu-das}{linkedin.com/in/apu-das}
\enspace | \enspace
\href{https://scholar.google.com/citations?user=GQbj1p4AAAAJ&hl=en}
{Google Scholar}

% =========================================================
% RESEARCH PROFILE
% =========================================================
\section{RESEARCH PROFILE}

PhD researcher studying hafnium-oxide-based ferroelectric
field-effect transistors (FeFETs), with a focus on memory
reliability, interface defects, and cryogenic operation.
Combines device fabrication and electrical characterization
from 300 K to 10 K to connect gate-stack physics with memory
performance. Research spans embedded memory, neuromorphic
devices, and industrial memory-development collaborations.
Interested in translating experimental insights into reliable,
energy-efficient semiconductor memory.

% =========================================================
% EDUCATION
% =========================================================
\section{EDUCATION}

\entry
  {Ph.D. Program, College of Semiconductor Research}
  {National Tsing Hua University (NTHU), Hsinchu, Taiwan}
  {2026 -- Present}

\begin{itemize}
  \item Converted from the M.Sc. program to the Ph.D. program
        in 2026; initially enrolled at NTHU in February 2025.
  \item Research advisor: Prof. Sourav De.
        Study focus: ferroelectric memory, interface engineering,
        device reliability, and cryogenic computing.
\end{itemize}

\entry
  {Bachelor of Technology, Electronics and Communication Engineering}
  {National Institute of Technology (NIT) Agartala, India}
  {June 2019 -- May 2023}

% =========================================================
% RESEARCH EXPERIENCE
% =========================================================
\section{RESEARCH EXPERIENCE}

\entry
  {Graduate Researcher -- Ferroelectric Memory Devices}
  {College of Semiconductor Research, NTHU, Taiwan}
  {February 2025 -- Present}

\begin{itemize}
  \item Characterize hafnium-oxide FeFETs from 300 K to 10 K
        to investigate switching behavior and reliability
        for embedded and cryogenic memory.
  \item Perform positive-up-negative-down (PUND) measurements
        directly on degraded transistors' gate stacks to
        compare polarization response with memory-window loss.
  \item Obtain evidence supporting interface-driven endurance
        degradation: polarization remains measurable after
        memory-window collapse near $10^{4}$--$10^{5}$ cycles
        in the studied devices.
  \item Use current--voltage (I--V), capacitance--voltage (C--V),
        and charge-pumping measurements to investigate charge
        trapping, oxygen-vacancy-related mechanisms, and
        cycling-induced threshold-voltage shifts.
  \item Study FeFET applications in spike-timing-dependent
        plasticity (STDP) and spiking neural networks.
\end{itemize}

\entry
  {Research Assistant -- VLSI}
  {Indian Institute of Technology (IIT) Guwahati, India}
  {2023 -- 2024}

% =========================================================
% COLLABORATIVE PROJECTS
% =========================================================
\section{INDUSTRIAL AND COLLABORATIVE PROJECTS}

% These are project affiliations, not claims of employment.

\textbf{Advanced Characterization and Reliability}\par
Taiwan Semiconductor Research Institute (TSRI) project |
2025 -- Present
\begin{itemize}
  \item Conduct endurance, retention, PUND, and charge-pumping
        measurements; develop pulse-measurement workflows and
        troubleshoot device behavior with fabrication,
        process, and equipment teams.
\end{itemize}

\textbf{2TnC DRAM Cell Development}\par
Powerchip Semiconductor (PSMC) collaboration | 2025
\begin{itemize}
  \item Contribute to cell design and process-flow definition,
        coordinating mask/layout alignment and deposition-related
        work with integration, lithography, and plasma-etch teams.
\end{itemize}

\textbf{FeRAM Chip Tapeout -- Ongoing}\par
Industrial Technology Research Institute (ITRI) project |
2025 -- Present
\begin{itemize}
  \item Collaborate on ferroelectric random-access memory
        (FeRAM) layout, integration, and tapeout preparation
        using TSRI EDA 2.0 and Cadence.
\end{itemize}

\newpage

% =========================================================
% PUBLICATIONS
% =========================================================
\section{SELECTED PUBLICATIONS (5 OF 14)}

% Verify publication status and equal-contribution designations
% against the final publication records before submission.

\textbf{Publication record:}
14 publications comprising 10 journal articles and
4 conference contributions. Selected first- and
joint-first-author journal articles are listed below.

\begin{enumerate}
  \item
  \textbf{``Sub-2 nm Equivalent-Oxide-Thickness Ferroelectric
  Transistors for Cryogenic Memory and Computing.''}\par
  \textit{ACS Nano}, 2026.
  \textbf{Joint first author.}

  \item
  \textbf{``Endurance Paradox in Hafnium-Oxide-Based
  Silicon-Channel Ferroelectric Transistors.''}\par
  \textit{ACS Applied Materials \& Interfaces}, 2026.
  \textbf{First author.}

  \item
  \textbf{``Vacancy-Assisted Polarization-Catastrophe Framework
  for Wake-Up and Ferroelectricity in Hafnium Oxide.''}\par
  \textit{Nano Letters}, 2026.
  \textbf{Joint first author.}

  \item
  \textbf{``Long-Term Reliability of Naturally Aged
  Hafnium-Oxide Ferroelectric Transistors for
  Energy-Efficient Embedded Memory.''}\par
  \textit{Cell Reports Physical Science}, 2026.
  \textbf{Joint first author.}

  \item
  \textbf{``Spike-Timing-Dependent Learning Dynamics in
  Silicon-Doped Hafnium-Oxide-Based Ferroelectric
  Field-Effect Transistors.''}\par
  \textit{IEEE Journal of the Electron Devices Society}, 2025.
  \textbf{Joint first author.}
\end{enumerate}

% =========================================================
% TECHNICAL SKILLS
% =========================================================
\section{TECHNICAL SKILLS}

\textbf{Devices and Reliability:}
FeFETs, FeRAM, DRAM, hafnium-oxide ferroelectrics,
high-k metal-gate (HKMG) stacks, endurance, retention,
charge trapping, interface defects, and root-cause analysis.\par

\textbf{Electrical Characterization:}
I--V, C--V, PUND, polarization--voltage (P--V),
charge pumping, pulse engineering, high-temperature testing,
and cryogenic measurements to 10 K.\par

\textbf{Instrumentation:}
Keysight B1500A with waveform generator/fast measurement unit
(WGFMU), Keithley 4200A, Radiant ferroelectric tester,
and cryogenic probe station.\par

\textbf{Fabrication and Integration:}
Process-flow definition, thin-film deposition,
lithography and plasma/dry-etch process tracking,
mask/layout alignment, and memory-device integration.\par

\textbf{EDA and Programming:}
Cadence Virtuoso, KLayout, TSRI EDA 2.0,
Python, Verilog, VHDL, Matplotlib, QtiPlot, and Linux.\par

\textbf{Cleanroom Qualification:}
TSRI SM01 and SM01-1 fabrication-facility operation
qualifications.\par

\textbf{Languages:}
English (fluent), Bengali (native), Mandarin (basic).

% =========================================================
% PATENT ACTIVITY
% =========================================================
\section{PATENT ACTIVITY}

\begin{itemize}
  \item \textbf{Taiwan patent application filed, March 2025:}
        ``Enhanced Independent Dual-Gate Ferroelectric
        Transistor Memory Using a Decoupled Ferroelectric
        Capacitor Structure.''
  \item \textbf{Patent filing in preparation, 2026:}
        ``2TnC Ferroelectric Memory Architecture Optimization,''
        in collaboration with PSMC.
\end{itemize}

% =========================================================
% TEACHING AND COMMUNICATION
% =========================================================
\section{TEACHING AND RESEARCH COMMUNICATION}

\begin{itemize}
  \item Served as a teaching assistant for two
        graduate/undergraduate courses.
  \item Presented research at the International Electron
        Devices and Materials Symposium (IEDMS) and
        Semiconductor Interface Specialists Conference (SISC).
\end{itemize}

% =========================================================
% SCHOLARSHIP MOTIVATION
% =========================================================
\section{SCHOLARSHIP MOTIVATION AND CAREER GOALS}

% Keep this personal statement only if it reflects your views.

My research has strengthened my interest in understanding why
memory devices degrade and how interface engineering can improve
their reliability. I am applying for the Micron Scholarship to
support my doctoral studies and connect my research interests
with industrial memory challenges. Having studied and conducted
research in India and Taiwan, I value international collaboration
and hope to contribute to an inclusive team through knowledge
sharing, respectful communication, and cross-disciplinary
problem-solving.

\end{document}
